\selectlanguage{ngerman}

\chapter*{Zusammenfassung}

Das in der Pampa Amarilla in Argentinien gelegene Pierre-Auger-Observatorium kombiniert verschiedene Nachweistechniken in einem Multi-Hybrid-Ansatz, um ultrahochenergetische kosmische Strahlung zu untersuchen.
Um die Messungen zu niedrigeren Energien zu erweitern, wurden innerhalb des 1500\,m-Arrays zwei dichtere Oberflächendetektor-Arrays errichtet: das \sd{750} und das \sd{433} mit Detektorabständen von 750\,m beziehungsweise 433\,m.
Diese Arbeit konzentriert sich auf das \sd{433}-Array, das mit Wasser-Cherenkov-Detektoren die Messung des Energiespektrums der kosmischen Strahlung ab $63\,\mathrm{PeV}$ ermöglicht.
Sein Energiebereich erlaubt die Charakterisierung der als zweites Knie bekannten Versteilung des Spektrums.
Eine vollständige Charakterisierung dieses Merkmals ist entscheidend für das Verständnis des Übergangs von galaktischer zu extragalaktischer kosmischer Strahlung.

Die Analyse basiert auf \SDFourThreeThreeNEvents{} Ereignissen, die zwischen Juli 2018 und Dezember 2021 mit Zenitwinkeln unter $40^\circ$ aufgezeichnet wurden.
Oberhalb der gewählten Energieschwelle kann das Array als vollständig effizient betrachtet werden, sodass sich die Exposition geometrisch berechnen lässt, ohne auf Simulationen zurückzugreifen.

Da nur wenige Ereignisse gleichzeitig vom \sd{433}-Array und den Fluoreszenzdetektoren gemessen werden, ist eine direkte Energiekalibrierung nicht möglich.
Ereignisse, die sowohl vom \sd{433} als auch vom \sd{750} gemessen werden, erlauben jedoch, die Schauergröße des \sd{433} anhand der vom \sd{750} bestimmten Energie zu kalibrieren.
Auf diese Weise wird die Energieskala des Fluoreszenzdetektors auf das \sd{433} übertragen.
Nach der Rekonstruktion der Primärenergie jedes Ereignisses wird das Rohspektrum bestimmt.
Zur Messung des zweiten Knies wird ein Likelihood-Fit durchgeführt, bei dem das Spektrummodell mit der Detektorantwort gefaltet wird.

Das Spektrum wird durch ein gebrochenes Potenzgesetz mit einem glatten Übergang beschrieben, der dem zweiten Knies entspricht.
Mit \sd{433}-Daten allein wird das zweite Knie bei $E_{01}=(2.8\pm0.5_{\mathrm{stat}}\pm0.7_{\mathrm{syst}})\times10^{17}\,\mathrm{eV}$ gemessen, wobei sich der Spektralindex von $\gamma_0=3.04\pm0.02_{\mathrm{stat}}\pm0.09_{\mathrm{syst}}$ unterhalb des Übergangs auf $\gamma_1=3.34\pm0.08_{\mathrm{stat}}\pm0.11_{\mathrm{syst}}$ oberhalb des Übergangs ändert.
Die Energieskala des Fluoreszenzdetektors dominiert die systematische Unsicherheit des Flusses und der Position des zweiten Knies, während die Kalibrierungsunsicherheit den größten Einfluss auf die Spektralindizes hat.

Außerdem wird eine Aktualisierung des 2021 veröffentlichten \sd{750}-Spektrums unter Verwendung der aktuellen Energieskala des Observatoriums vorgestellt.
Mit den von beiden Arrays aufgezeichneten Ereignissen wird auch die Nachweiseffizienz des \sd{750} neu bestimmt.
Da die Position des zweiten Knies mit \sd{750}-Daten allein nicht gemessen werden kann, wird aus den Daten beider Arrays ein vorläufiges kombiniertes Spektrum bestimmt.
Ziel dieser Kombination ist es, die Position des zweiten Knies besser einzugrenzen und die Spektrummessung bis über den Knöchel hinaus zu erweitern, um so zum Verständnis des Übergangs von galaktischer zu extragalaktischer kosmischer Strahlung beizutragen.
